% from main.tex
%%%%%%%%%%  scan idx 73  |  folio 60  %%%%%%%%%%
Utilisant le graphe de $f$, on décompose $f$ en produit d'une \emph{immersion} de $X$ dans $X\times Y$ et d'une \emph{projection} de $X\times Y$ sur $Y$. Pour cette dernière le corollaire est trivial, et pour une immersion, on procède comme pour le corollaire 1, en utilisant la proposition 2.8 à la place de la proposition 2.7.

\textbf{Proposition 2.13 (``décomposition cellulaire'').} \emph{Soit $P$ un espace algébrique (irréductible ou non), muni d'une famille croissante $P^0\subset P^1\subset\cdots\subset P^n=P$ de parties fermées, telles que pour tout $i$ avec $0\leq i\leq n$, les composantes connexes de $P^i-P^{i-1}$ soient des espaces algébriques $V_{i,\alpha}$ isomorphes à $k^i$ (on pose $P^{-1}=\varnothing$). Alors~:}
\begin{enumerate}[label=\alph*)]
\item \emph{Pour tout espace algébrique $X$, l'application naturelle de $\underline K(X)\otimes\underline K(P)$ dans $\underline K(X\times P)$ déduite du produit tensoriel de faisceaux sur $X$ et $P$, est surjective.}
\item \emph{Les $\gamma_P(V_{i,\alpha})$ forment un système de générateurs du groupe $\underline K(P)$.}
\end{enumerate}
La démonstration se fait par récurrence sur $n$, en utilisant le corollaire de la proposition 2.9 et la proposition 2.10. (N.B. Malheureusement, j'ignore si les $\gamma_P(V_{i,\alpha})$ forment une base de $\underline K(P)$ et si l'homomorphisme envisagé dans a) est bijectif. Nous n'aurons pas besoin de le savoir pour la suite).

\textbf{Théorème 2.14.}
